\documentclass[12pt,a4paper]{article}
\usepackage[T1]{fontenc}
\usepackage{amsmath,amssymb,mathrsfs,tikz,times,pifont}
\usepackage{enumitem}
\usepackage{multicol}
\usepackage{lmodern}
\usetikzlibrary{trees}
\newcommand\circitem[1]{%
\tikz[baseline=(char.base)]{
\node[circle,draw=gray, fill=red!55,
minimum size=1.2em,inner sep=0] (char) {#1};}}
\newcommand\boxitem[1]{%
\tikz[baseline=(char.base)]{
\node[fill=cyan,
minimum size=1.2em,inner sep=0] (char) {#1};}}
\setlist[enumerate,1]{label=\protect\circitem{\arabic*}}
\setlist[enumerate,2]{label=\protect\boxitem{\alph*}}
%%%::::::by chnini ameur :::::::%%%
\everymath{\displaystyle}
\usepackage[left=1cm,right=1cm,top=1cm,bottom=1.7cm]{geometry}
\usepackage[colorlinks=true, linkcolor=blue, urlcolor=blue, citecolor=blue]{hyperref}
\usepackage{array,multirow}
\usepackage[most]{tcolorbox}
\usepackage{varwidth}
\usepackage{float} %pour utiliser l'option [H] qui force l'image à apparaître exactement à l'endroit où elle est placée dans le code.
\tcbuselibrary{skins,hooks}
\usetikzlibrary{patterns}
%%%::::::by chnini ameur :::::::%%%
\newtcolorbox{exa}[2][]{enhanced,breakable,before skip=2mm,after skip=5mm,
colback=yellow!20!white,colframe=black!20!blue,boxrule=0.5mm,
attach boxed title to top left ={xshift=0.6cm,yshift*=1mm-\tcboxedtitleheight},
fonttitle=\bfseries,
title={#2},#1,
% varwidth boxed title*=-3cm,
boxed title style={frame code={
\path[fill=tcbcolback!30!black]
([yshift=-1mm,xshift=-1mm]frame.north west)
arc[start angle=0,end angle=180,radius=1mm]
([yshift=-1mm,xshift=1mm]frame.north east)
arc[start angle=180,end angle=0,radius=1mm];
\path[left color=tcbcolback!60!black,right color = tcbcolback!60!black,
middle color = tcbcolback!80!black]
([xshift=-2mm]frame.north west) -- ([xshift=2mm]frame.north east)
[rounded corners=1mm]-- ([xshift=1mm,yshift=-1mm]frame.north east)
-- (frame.south east) -- (frame.south west)
-- ([xshift=-1mm,yshift=-1mm]frame.north west)
[sharp corners]-- cycle;
},interior engine=empty,
},interior style={top color=yellow!5}}
%%%%%%%%%%%%%%%%%%%%%%%
\usepackage{fancyhdr}
\usepackage{eso-pic}         % Pour ajouter des éléments en arrière-plan
% Commande pour ajouter du texte en arrière-plan
\usepackage{tkz-tab}
\usepackage{tikz}
\usepackage[siunitx, european, RPvoltages]{circuitikz}
\AddToShipoutPicture{
    \AtTextCenter{%
        \makebox[0pt]{\rotatebox{80}{\textcolor[gray]{0.7}{\fontsize{5cm}{5cm}\selectfont PGB}}}
    }
}
\usepackage{lastpage}
\fancyhf{}
\pagestyle{fancy}
\renewcommand{\footrulewidth}{1pt}
\renewcommand{\headrulewidth}{0pt}
\renewcommand{\footruleskip}{10pt}
\fancyfoot[R]{
\color{blue}\ding{45}\ \textbf{2026}
}
\fancyfoot[L]{
\color{blue}\ding{45}\ \textbf{Prof:M. BA}
}
\cfoot{\bf
\thepage /
\pageref{LastPage}}
% Création du compteur pour les exercices
\newcounter{exercice}
\renewcommand{\theexercice}{\arabic{exercice}}  % Définit l'affichage du compteur en chiffres arabes

% Définir la commande \exo
\newcommand{\exo}{\refstepcounter{exercice}\textbf{Exercice \theexercice} }
\begin{document}
\renewcommand{\arraystretch}{1.5}
\renewcommand{\arrayrulewidth}{1.2pt}
\begin{tikzpicture}[overlay,remember picture]
    \node[draw=blue,line width=1.2pt,fill=purple,text=blue,inner sep=3mm,rounded corners,pattern=dots]at ([yshift=-2.5cm]current page.north) {\begingroup\setlength{\fboxsep}{0pt}\colorbox{white}{\begin{tabular}{|*1{>{\centering \arraybackslash}p{0.28\textwidth}} |*2{>{\centering \arraybackslash}p{0.2\textwidth}|} *1{>{\centering \arraybackslash}p{0.19\textwidth}|} }
                \hline
                \multicolumn{3}{|c|}{$\diamond$$\diamond$$\diamond$\ \textbf{Lycée de Dindéfélo}\ $\diamond$$\diamond$$\diamond$ } & \textbf{A.S. : 2026/2027}                                              \\ \hline
                \textbf{Matière: Mathématiques}                                                                                    & \textbf{Niveau : 3}\textbf{M1} & \textbf{Date: 17/11/2026} & \textbf{} \\ \hline
                \multicolumn{4}{|c|}{\parbox[c]{10cm}{\begin{center}
                                                                  \textbf{{\Large\sffamily Correction Td Racine Carrée}}
                                                              \end{center}}}                                                                                                        \\ \hline
            \end{tabular}}\endgroup};
\end{tikzpicture}

\vspace{3cm}
\fbox{\textbf{\exo}}
Écrire sous la forme $a\sqrt{b}$ où $a$ et $b$ sont des nombres entiers, $b$ étant le plus petit possible :
\begin{multicols}{2}
\setlength{\columnseprule}{0.1mm} % La largeur de la ligne verticale entre les colonnes
\begin{enumerate}
    \item \textbf{Calcul de $A$ :}
    
    $
    \begin{aligned}
    A &= \sqrt{50} \\
    &= \sqrt{5 \times 5 \times 2} \\
    &= \sqrt{5^2 \times 2} \\
    &= \sqrt{5^2} \times \sqrt{2} \\
    &= 5 \times \sqrt{2} \\
    &= 5\sqrt{2}
    \end{aligned}
    $

    \item \textbf{Calcul de $B$ :}
    
    $
    \begin{aligned}
    B &= \sqrt{72} \\
    &= \sqrt{2 \times 2 \times 2 \times 3 \times 3} \\
    &= \sqrt{2^2 \times 3^2 \times 2} \\
    &= \sqrt{2^2} \times \sqrt{3^2} \times \sqrt{2} \\
    &= 2 \times 3 \times \sqrt{2} \\
    &= 6\sqrt{2}
    \end{aligned}
    $

    \item \textbf{Calcul de $C$ :}
    
    $
    \begin{aligned}
    C &= \sqrt{50} + \sqrt{72} \\
    &= \sqrt{5^2 \times 2} + \sqrt{6^2 \times 2} \\
    &= \sqrt{5^2} \times \sqrt{2} + \sqrt{6^2} \times \sqrt{2} \\
    &= 5\sqrt{2} + 6\sqrt{2} \\
    &= (5 + 6)\sqrt{2} \\
    &= 11\sqrt{2}
    \end{aligned}
    $

    \item \textbf{Calcul de $D$ :}
    
    $
    \begin{aligned}
    D &= 2\sqrt{3} + \sqrt{75} - 6\sqrt{27} \\
    &= 2\sqrt{3} + \sqrt{5^2 \times 3} - 6\sqrt{3^2 \times 3} \\
    &= 2\sqrt{3} + \sqrt{5^2}\times\sqrt{3} - 6\sqrt{3^2}\times\sqrt{3} \\
    &= 2\sqrt{3} + 5\sqrt{3} - 6 \times 3 \times \sqrt{3} \\
    &= 2\sqrt{3} + 5\sqrt{3} - 18\sqrt{3} \\
    &= (2 + 5 - 18)\sqrt{3} \\
    &= -11\sqrt{3}
    \end{aligned}
    $

    \item \textbf{Calcul de $E$ :}
    
    $
    \begin{aligned}
    E &= 2\sqrt{3} \times \sqrt{6} \\
    &= 2\sqrt{3 \times (2 \times 3)} \\
    &= 2\sqrt{3^2 \times 2} \\
    &= 2 \times \sqrt{3^2} \times \sqrt{2} \\
    &= 2 \times 3 \times \sqrt{2} \\
    &= 6\sqrt{2}
    \end{aligned}
    $

    \item \textbf{Calcul de $F$ :}
    
    $
    \begin{aligned}
    F &= \sqrt{8} \times \sqrt{50} \times \sqrt{18} \\
    &= \sqrt{2^3} \times \sqrt{5^2 \times 2} \times \sqrt{3^2 \times 2} \\
    &= \sqrt{2^2 \times 2} \times (5\sqrt{2}) \times (3\sqrt{2}) \\
    &= 2\sqrt{2} \times 5\sqrt{2} \times 3\sqrt{2} \\
    &= (2 \times 5 \times 3) \times (\sqrt{2} \times \sqrt{2} \times \sqrt{2}) \\
    &= 30 \times (\sqrt{2})^2 \times \sqrt{2} \\
    &= 30 \times 2 \times \sqrt{2} \\
    &= 60\sqrt{2}
    \end{aligned}
    $
\end{enumerate}
\end{multicols}
\newpage
\fbox{\textbf{\exo}}

\begin{enumerate}
    \item \textbf{Calcul de $A$ :}
    
    $
    \begin{aligned}
    A &= \sqrt{49} + \sqrt{28} + \sqrt{63} \\
    &= \sqrt{7 \times 7} + \sqrt{2 \times 2 \times 7} + \sqrt{3 \times 3 \times 7} \\
    &= \sqrt{7^2} + \sqrt{2^2 \times 7} + \sqrt{3^2 \times 7} \\
    &= 7 + \sqrt{2^2} \times \sqrt{7} + \sqrt{3^2} \times \sqrt{7} \\
    &= 7 + 2\sqrt{7} + 3\sqrt{7} \\
    &= 7 + (2 + 3)\sqrt{7} \\
    &= 7 + 5\sqrt{7} \textbf{ \textcolor{blue}{ Par Aïssatou Sadio Sylla}}
    \end{aligned}
    $

    \item \textbf{Calcul de $B$ :}
    
    $
    \begin{aligned}
    B &= (2\sqrt{7} + 1)^2 - (\sqrt{3} - 1)(\sqrt{3} + 1) \\
    &= \left[(2\sqrt{7})^2 + 2 \times 2\sqrt{7} \times 1 + 1^2\right] - \left[(\sqrt{3})^2 - 1^2\right] \\
    &= \left(4 \times 7 + 4\sqrt{7} + 1\right) - (3 - 1) \\
    &= (28 + 4\sqrt{7} + 1) - 2 \\
    &= 29 + 4\sqrt{7} - 2 \\
    &= 27 + 4\sqrt{7} \textbf{ \textcolor{blue}{ Par Goulo Touré}}
    \end{aligned}
    $

    \item \textbf{Calcul de $C$ :}
    
    $
    \begin{aligned}
    C &= 6\sqrt{28} + 10\sqrt{7} - 8\sqrt{63} \\
    &= 6\sqrt{2 \times 2 \times 7} + 10\sqrt{7} - 8\sqrt{3 \times 3 \times 7} \\
    &= 6\sqrt{2^2 \times 7} + 10\sqrt{7} - 8\sqrt{3^2 \times 7} \\
    &= 6 \times \sqrt{2^2} \times \sqrt{7} + 10\sqrt{7} - 8 \times \sqrt{3^2} \times \sqrt{7} \\
    &= 6 \times 2\sqrt{7} + 10\sqrt{7} - 8 \times 3\sqrt{7} \\
    &= 12\sqrt{7} + 10\sqrt{7} - 24\sqrt{7} \\
    &= (12 + 10 - 24)\sqrt{7} \\
    &= -2\sqrt{7} \\
    &= 0 - 2\sqrt{7} \textbf{ \textcolor{blue}{ Par Aïssatou Diallo }}
    \end{aligned}
    $
\end{enumerate}

\fbox{\textbf{\exo}}

\begin{itemize}
    \item \textbf{Calcul de $D$ :} \\
    $
    \begin{aligned}
    D &= (2\sqrt{3} + 1) \times (2\sqrt{3} - 1) \\
    &= (2\sqrt{3})^2 - 1^2 \\
    &= 2^2 \times (\sqrt{3})^2 - 1 \\
    &= 4 \times 3 - 1 \\
    &= 12 - 1 \\
    &= 11 \textbf{ \textcolor{blue}{ Par ............}}
    \end{aligned}
    $

    \item \textbf{Calcul de $E$ :} \\
    $
    \begin{aligned}
    E &= 8\sqrt{5} - \sqrt{20} - 2\sqrt{45} \\
    &= 8\sqrt{5} - \sqrt{2 \times 2 \times 5} - 2\sqrt{3 \times 3 \times 5} \\
    &= 8\sqrt{5} - \sqrt{2^2 \times 5} - 2\sqrt{3^2 \times 5} \\
    &= 8\sqrt{5} - \sqrt{2^2} \times \sqrt{5} - 2 \times \sqrt{3^2} \times \sqrt{5} \\
    &= 8\sqrt{5} - 2\sqrt{5} - 2 \times 3\sqrt{5} \\
    &= 8\sqrt{5} - 2\sqrt{5} - 6\sqrt{5} \\
    &= (8 - 2 - 6)\sqrt{5} \\
    &= 0\sqrt{5} \\
    &= 0 \textbf{ \textcolor{blue}{ Par ..........}}
    \end{aligned}
    $
\end{itemize}


\fbox{\textbf{\exo}}

\begin{enumerate}
    \item \textbf{Écrivons $\sqrt{5} \times \sqrt{125}$ sous la forme d'un nombre entier :} \\
    $
    \begin{aligned}
    \sqrt{5} \times \sqrt{125} &= \sqrt{5} \times \sqrt{5 \times 5 \times 5} \\
    &= \sqrt{5} \times \sqrt{5^2 \times 5} \\
    &= \sqrt{5} \times \sqrt{5^2} \times \sqrt{5} \\
    &= \sqrt{5} \times 5 \times \sqrt{5} \\
    &= 5 \times (\sqrt{5})^2 \\
    &= 5 \times 5 \\
    &= 25 \textbf{ \textcolor{blue}{ Par ..........}}
    \end{aligned}
    $

    \item \textbf{Écrivons $(\sqrt{5} \times \sqrt{125}) \times 2\sqrt{5}$ sous la forme $a\sqrt{5}$ :} \\
    D'après la question précédente, $\sqrt{5} \times \sqrt{125} = 25$, donc : \\
    $
    \begin{aligned}
    (\sqrt{5} \times \sqrt{125}) \times 2\sqrt{5} &= 25 \times 2\sqrt{5} \\
    &= (25 \times 2)\sqrt{5} \\
    &= 50\sqrt{5} \textbf{ \textcolor{blue}{ Par ..........}}
    \end{aligned}
    $
\end{enumerate}

\fbox{\textbf{\exo}}

\begin{enumerate}
    \item Écrivons $A$, $B$ et $C$ sous la forme $a\sqrt{3}$ :
    \begin{enumerate}
        \item \textbf{Calcul de $A$ :} \\
        $
        \begin{aligned}
        A &= \sqrt{12} + 5\sqrt{75} - 2\sqrt{27} \\
        &= \sqrt{2 \times 2 \times 3} + 5\sqrt{5 \times 5 \times 3} - 2\sqrt{3 \times 3 \times 3} \\
        &= \sqrt{2^2 \times 3} + 5\sqrt{5^2 \times 3} - 2\sqrt{3^2 \times 3} \\
        &= \sqrt{2^2} \times \sqrt{3} + 5 \times \sqrt{5^2} \times \sqrt{3} - 2 \times \sqrt{3^2} \times \sqrt{3} \\
        &= 2\sqrt{3} + 5 \times 5\sqrt{3} - 2 \times 3\sqrt{3} \\
        &= 2\sqrt{3} + 25\sqrt{3} - 6\sqrt{3} \\
        &= (2 + 25 - 6)\sqrt{3} \\
        &= 21\sqrt{3} \textbf{ \textcolor{blue}{ Par Karim Baldé}}
        \end{aligned}
        $

        \item \textbf{Calcul de $B$ :} \\
        
        $
        \begin{aligned}
        B &= (5 + \sqrt{3})^2 - (2\sqrt{7})^2 \\
        &= \left[5^2 + 2 \times 5 \times \sqrt{3} + (\sqrt{3})^2\right] - \left[2^2 \times (\sqrt{7})^2\right] \\
        &= (25 + 10\sqrt{3} + 3) - (4 \times 7) \\
        &= 28 + 10\sqrt{3} - 28 \\
        &= 10\sqrt{3} \textbf{ \textcolor{blue}{ Par Aïssatou Sadio Sylla}}
        \end{aligned}
        $
		
        \item \textbf{Calcul de $C$ :} \\
        $
        \begin{aligned}
        C &= \sqrt{12} - \sqrt{75} - 2\sqrt{27} \\
        &= \sqrt{2 \times 2 \times 3} - \sqrt{5 \times 5 \times 3} - 2\sqrt{3 \times 3 \times 3} \\
        &= \sqrt{2^2 \times 3} - \sqrt{5^2 \times 3} - 2\sqrt{3^2 \times 3} \\
        &= \sqrt{2^2} \times \sqrt{3} - \sqrt{5^2} \times \sqrt{3} - 2 \times \sqrt{3^2} \times \sqrt{3} \\
        &= 2\sqrt{3} - 5\sqrt{3} - 2 \times 3\sqrt{3} \\
        &= 2\sqrt{3} - 5\sqrt{3} - 6\sqrt{3} \\
        &= (2 - 5 - 6)\sqrt{3} \\
        &= -9\sqrt{3} \textbf{ \textcolor{blue}{ Par Aïssatou Diallo}}
        \end{aligned}
        $
    \end{enumerate}

    \item \textbf{Produit des nombres $E$ et $F$ :} \\
        
    $
    \begin{aligned}
    E \times F &= (3\sqrt{5} - 2\sqrt{11})(3\sqrt{5} + 2\sqrt{11}) \\
    &= (3\sqrt{5})^2 - (2\sqrt{11})^2 \\
    &= 3^2 \times (\sqrt{5})^2 - 2^2 \times (\sqrt{11})^2 \\
    &= 9 \times 5 - 4 \times 11 \\
    &= 45 - 44 \\
    &= 1
    \end{aligned}
    $
    		
    		Donc E et F sont des inverses \textbf{ \textcolor{blue}{ Par Goulo Touré}}
		\end{enumerate}

\fbox{\textbf{\exo}}

\begin{comment}

\begin{enumerate}
    \item \begin{enumerate}
        \item \textbf{Écrivons $A$ sous la forme $a\sqrt{3} + b\sqrt{5}$ :} \\
        $
        \begin{aligned}
        A &= \sqrt{48} + \sqrt{20} \\
        &= \sqrt{2 \times 2 \times 2 \times 2 \times 3} + \sqrt{2 \times 2 \times 5} \\
        &= \sqrt{4^2 \times 3} + \sqrt{2^2 \times 5} \\
        &= \sqrt{4^2} \times \sqrt{3} + \sqrt{2^2} \times \sqrt{5} \\
        &= 4\sqrt{3} + 2\sqrt{5}
        \end{aligned}
        $ \\
        Donc $A$ s'écrit bien sous la forme $a\sqrt{3} + b\sqrt{5}$ avec $a = 4$ et $b = 2$.

        \item \textbf{Écrivons $B$ sous la forme $c\sqrt{3} + d\sqrt{5}$ :} \\
        $
        \begin{aligned}
        B &= \sqrt{108} - \sqrt{45} \\
        &= \sqrt{2 \times 2 \times 3 \times 3 \times 3} - \sqrt{3 \times 3 \times 5} \\
        &= \sqrt{6^2 \times 3} - \sqrt{3^2 \times 5} \\
        &= \sqrt{6^2} \times \sqrt{3} - \sqrt{3^2} \times \sqrt{5} \\
        &= 6\sqrt{3} - 3\sqrt{5}
        \end{aligned}
        $ \\
        Donc $B$ s'écrit bien sous la forme $c\sqrt{3} + d\sqrt{5}$ avec $c = 6$ et $d = -3$.
    \end{enumerate}

    \item \textbf{Montrons que le produit $AB$ est un nombre entier :} \\
    $
    \begin{aligned}
    AB &= (4\sqrt{3} + 2\sqrt{5})(6\sqrt{3} - 3\sqrt{5}) \\
    &= 2(2\sqrt{3} + \sqrt{5}) \times 3(2\sqrt{3} - \sqrt{5}) \\
    &= 6(2\sqrt{3} + \sqrt{5})(2\sqrt{3} - \sqrt{5}) \\
    &= 6\left[(2\sqrt{3})^2 - (\sqrt{5})^2\right] \\
    &= 6\left[2^2 \times (\sqrt{3})^2 - 5\right] \\
    &= 6(4 \times 3 - 5) \\
    &= 6(12 - 5) \\
    &= 6 \times 7 \\
    &= 42
    \end{aligned}
    $ \\
    $42$ est un nombre entier, donc $AB$ est bien un nombre entier.
\end{enumerate}

\end{comment}

\fbox{\textbf{\exo}}

\begin{comment}

Simplifions d'abord l'écriture de $A$ : \\
$
\begin{aligned}
A &= \sqrt{27} + 1 \\
&= \sqrt{3 \times 3 \times 3} + 1 \\
&= \sqrt{3^2 \times 3} + 1 \\
&= \sqrt{3^2} \times \sqrt{3} + 1 \\
&= 3\sqrt{3} + 1
\end{aligned}
$

\begin{itemize}
    \item \textbf{Calcul de $A - B$ :} \\
    $
    \begin{aligned}
    A - B &= (3\sqrt{3} + 1) - (2\sqrt{3} - 5) \\
    &= 3\sqrt{3} + 1 - 2\sqrt{3} + 5 \\
    &= (3 - 2)\sqrt{3} + (1 + 5) \\
    &= 1\sqrt{3} + 6 \\
    &= \sqrt{3} + 6
    \end{aligned}
    $

    \item \textbf{Calcul de $A^2$ :} \\
    $
    \begin{aligned}
    A^2 &= (3\sqrt{3} + 1)^2 \\
    &= (3\sqrt{3})^2 + 2 \times 3\sqrt{3} \times 1 + 1^2 \\
    &= 3^2 \times (\sqrt{3})^2 + 6\sqrt{3} + 1 \\
    &= 9 \times 3 + 6\sqrt{3} + 1 \\
    &= 27 + 6\sqrt{3} + 1 \\
    &= 6\sqrt{3} + 28
    \end{aligned}
    $

    \item \textbf{Calcul de $B^2$ :} \\
    $
    \begin{aligned}
    B^2 &= (2\sqrt{3} - 5)^2 \\
    &= (2\sqrt{3})^2 - 2 \times 2\sqrt{3} \times 5 + 5^2 \\
    &= 2^2 \times (\sqrt{3})^2 - 20\sqrt{3} + 25 \\
    &= 4 \times 3 - 20\sqrt{3} + 25 \\
    &= 12 - 20\sqrt{3} + 25 \\
    &= -20\sqrt{3} + 37
    \end{aligned}
    $
\end{itemize}

\end{comment}

\fbox{\textbf{\exo}}

\begin{comment}

\begin{itemize}
    \item \textbf{Calcul de $D - E$ :} \\
    $
    \begin{aligned}
    D - E &= (5 - 3\sqrt{2}) - (4 + 5\sqrt{2}) \\
    &= 5 - 3\sqrt{2} - 4 - 5\sqrt{2} \\
    &= (5 - 4) + (-3 - 5)\sqrt{2} \\
    &= 1 - 8\sqrt{2}
    \end{aligned}
    $

    \item \textbf{Calcul de $D \times E$ :} \\
    $
    \begin{aligned}
    D \times E &= (5 - 3\sqrt{2})(4 + 5\sqrt{2}) \\
    &= 5 \times 4 + 5 \times 5\sqrt{2} - 3\sqrt{2} \times 4 - 3\sqrt{2} \times 5\sqrt{2} \\
    &= 20 + 25\sqrt{2} - 12\sqrt{2} - 15 \times (\sqrt{2})^2 \\
    &= 20 + (25 - 12)\sqrt{2} - 15 \times 2 \\
    &= 20 + 13\sqrt{2} - 30 \\
    &= -10 + 13\sqrt{2}
    \end{aligned}
    $
\end{itemize}

\end{comment}

\fbox{\textbf{\exo}}

\begin{comment}

\begin{enumerate}
    \item \textbf{Calcul de $C$ :} \\
    $
    \begin{aligned}
    C &= \sqrt{500} + 3\sqrt{5} - 3\sqrt{45} \\
    &= \sqrt{2 \times 2 \times 5 \times 5 \times 5} + 3\sqrt{5} - 3\sqrt{3 \times 3 \times 5} \\
    &= \sqrt{10^2 \times 5} + 3\sqrt{5} - 3\sqrt{3^2 \times 5} \\
    &= \sqrt{10^2} \times \sqrt{5} + 3\sqrt{5} - 3 \times \sqrt{3^2} \times \sqrt{5} \\
    &= 10\sqrt{5} + 3\sqrt{5} - 3 \times 3\sqrt{5} \\
    &= 10\sqrt{5} + 3\sqrt{5} - 9\sqrt{5} \\
    &= (10 + 3 - 9)\sqrt{5} \\
    &= 4\sqrt{5}
    \end{aligned}
    $

    \item \textbf{Calcul de $D$ :} \\
    $
    \begin{aligned}
    D &= (5 - 2\sqrt{6})(5 + 2\sqrt{6}) \\
    &= 5^2 - (2\sqrt{6})^2 \\
    &= 25 - 2^2 \times (\sqrt{6})^2 \\
    &= 25 - 4 \times 6 \\
    &= 25 - 24 \\
    &= 1
    \end{aligned}
    $
\end{enumerate}

\end{comment}



\fbox{\textbf{\exo}}

\begin{comment}

\begin{itemize}
    \item \textbf{Calcul de $B$ :} \\
    $
    \begin{aligned}
    B &= (5 + \sqrt{3})(5 - \sqrt{3}) - 8\sqrt{5}(\sqrt{5} - 1) \\
    &= \left[5^2 - (\sqrt{3})^2\right] - \left[8(\sqrt{5})^2 - 8\sqrt{5} \times 1\right] \\
    &= (25 - 3) - (8 \times 5 - 8\sqrt{5}) \\
    &= 22 - (40 - 8\sqrt{5}) \\
    &= 22 - 40 + 8\sqrt{5} \\
    &= -18 + 8\sqrt{5}
    \end{aligned}
    $
\end{itemize}

\end{comment}



\fbox{\textbf{\exo}}

\begin{comment}

\begin{enumerate}
    \item \textbf{Calcul de $a^2$, $b^2$ et $a^2 + b^2$ :}
    \begin{itemize}
        \item \textbf{Calcul de $a^2$ :} \\
        $
        \begin{aligned}
        a^2 &= \left[\sqrt{3}(1 + \sqrt{6})\right]^2 \\
        &= (\sqrt{3})^2 \times (1 + \sqrt{6})^2 \\
        &= 3 \times \left[1^2 + 2 \times 1 \times \sqrt{6} + (\sqrt{6})^2\right] \\
        &= 3 \times (1 + 2\sqrt{6} + 6) \\
        &= 3 \times (7 + 2\sqrt{6}) \\
        &= 21 + 6\sqrt{6}
        \end{aligned}
        $

        \item \textbf{Calcul de $b^2$ :} \\
        $
        \begin{aligned}
        b^2 &= (3 - \sqrt{6})^2 \\
        &= 3^2 - 2 \times 3 \times \sqrt{6} + (\sqrt{6})^2 \\
        &= 9 - 6\sqrt{6} + 6 \\
        &= 15 - 6\sqrt{6}
        \end{aligned}
        $

        \item \textbf{Calcul de $a^2 + b^2$ :} \\
        $
        \begin{aligned}
        a^2 + b^2 &= (21 + 6\sqrt{6}) + (15 - 6\sqrt{6}) \\
        &= 21 + 15 + 6\sqrt{6} - 6\sqrt{6} \\
        &= 36
        \end{aligned}
        $
    \end{itemize}

    \item \textbf{Nature de $a^2 + b^2$ :} \\
    On a trouvé $a^2 + b^2 = 36$. Comme $36 \in \mathbb{N}$, $a^2 + b^2$ est bien un nombre entier.

    \item \textbf{Longueur de l'hypoténuse :} \\
    Soit $c$ la longueur de l'hypoténuse ($c > 0$). D'après le théorème de Pythagore : \\
    $
    \begin{aligned}
    c^2 &= a^2 + b^2 \\
    c^2 &= 36 \\
    c &= \sqrt{36} \quad (\text{car } c > 0) \\
    c &= 6
    \end{aligned}
    $ \\
    La longueur de l'hypoténuse est égale à $6$.
\end{enumerate}

\end{comment}

\fbox{\textbf{\exo}}

\begin{enumerate}
    \item \textbf{Simplification de $A$ :} \\
    $
    \begin{aligned}
    A &= \dfrac{\sqrt{8} + 2\sqrt{32}}{5\sqrt{2} + \sqrt{50}} \\[0.2cm]
    &= \dfrac{\sqrt{2 \times 2 \times 2} + 2\sqrt{4 \times 4 \times 2}}{5\sqrt{2} + \sqrt{5 \times 5 \times 2}} \\[0.2cm]
    &= \dfrac{\sqrt{2^2 \times 2} + 2\sqrt{4^2 \times 2}}{5\sqrt{2} + \sqrt{5^2 \times 2}} \\[0.2cm]
    &= \dfrac{\sqrt{2^2} \times \sqrt{2} + 2 \times \sqrt{4^2} \times \sqrt{2}}{5\sqrt{2} + \sqrt{5^2} \times \sqrt{2}} \\[0.2cm]
    &= \dfrac{2\sqrt{2} + 2 \times 4\sqrt{2}}{5\sqrt{2} + 5\sqrt{2}} \\[0.2cm]
    &= \dfrac{2\sqrt{2} + 8\sqrt{2}}{10\sqrt{2}} \\[0.2cm]
    &= \dfrac{10\sqrt{2}}{10\sqrt{2}} \\[0.2cm]
    &= 1
    \end{aligned}
    $

    \item \textbf{Simplification de $B$ :} \\
    $
    \begin{aligned}
    B &= \sqrt{2 - \sqrt{3}} \times \sqrt{2 + \sqrt{3}} \\
    &= \sqrt{(2 - \sqrt{3})(2 + \sqrt{3})} \\
    &= \sqrt{2^2 - (\sqrt{3})^2} \\
    &= \sqrt{4 - 3} \\
    &= \sqrt{1} \\
    &= 1
    \end{aligned}
    $

    \item \textbf{Simplification de $C$ :} \\
    Rendons rationnels les dénominateurs de chaque fraction en multipliant par leurs expressions conjuguées : \\
    $
    \begin{aligned}
    C &= \dfrac{\sqrt{2} + 1}{\sqrt{2} - 1} + \dfrac{2\sqrt{2} + 1}{\sqrt{2} + 1} \\[0.2cm]
    &= \dfrac{(\sqrt{2} + 1)(\sqrt{2} + 1)}{(\sqrt{2} - 1)(\sqrt{2} + 1)} + \dfrac{(2\sqrt{2} + 1)(\sqrt{2} - 1)}{(\sqrt{2} + 1)(\sqrt{2} - 1)} \\[0.2cm]
    &= \dfrac{(\sqrt{2} + 1)^2}{(\sqrt{2})^2 - 1^2} + \dfrac{2(\sqrt{2})^2 - 2\sqrt{2} + \sqrt{2} - 1}{(\sqrt{2})^2 - 1^2} \\[0.2cm]
    &= \dfrac{(\sqrt{2})^2 + 2\sqrt{2} + 1^2}{2 - 1} + \dfrac{2 \times 2 - \sqrt{2} - 1}{2 - 1} \\[0.2cm]
    &= \dfrac{2 + 2\sqrt{2} + 1}{1} + \dfrac{4 - \sqrt{2} - 1}{1} \\[0.2cm]
    &= (3 + 2\sqrt{2}) + (3 - \sqrt{2}) \\
    &= (3 + 3) + (2 - 1)\sqrt{2} \\
    &= 6 + \sqrt{2}
    \end{aligned}
    $
\end{enumerate}

\fbox{\textbf{\exo}}

\begin{comment}
\noindent\textbf{\textit{Correction :}} \\[0.2cm]

\begin{enumerate}
    \item \begin{itemize}
        \item \textbf{Rendons rationnel le dénominateur de $a$ :} \\
        $
        \begin{aligned}
        a &= \dfrac{2 - \sqrt{3}}{5 + \sqrt{3}} \\[0.2cm]
        &= \dfrac{(2 - \sqrt{3})(5 - \sqrt{3})}{(5 + \sqrt{3})(5 - \sqrt{3})} \\[0.2cm]
        &= \dfrac{2 \times 5 - 2\sqrt{3} - 5\sqrt{3} + (\sqrt{3})^2}{5^2 - (\sqrt{3})^2} \\[0.2cm]
        &= \dfrac{10 - 7\sqrt{3} + 3}{25 - 3} \\[0.2cm]
        &= \dfrac{13 - 7\sqrt{3}}{22}
        \end{aligned}
        $

        \item \textbf{Simplifions $b$ :} \\
        $
        \begin{aligned}
        b &= 3\sqrt{18} + \sqrt{128} - \sqrt{338} \\
        &= 3\sqrt{3 \times 3 \times 2} + \sqrt{8 \times 8 \times 2} - \sqrt{13 \times 13 \times 2} \\
        &= 3\sqrt{3^2 \times 2} + \sqrt{8^2 \times 2} - \sqrt{13^2 \times 2} \\
        &= 3 \times \sqrt{3^2} \times \sqrt{2} + \sqrt{8^2} \times \sqrt{2} - \sqrt{13^2} \times \sqrt{2} \\
        &= 3 \times 3\sqrt{2} + 8\sqrt{2} - 13\sqrt{2} \\
        &= 9\sqrt{2} + 8\sqrt{2} - 13\sqrt{2} \\
        &= (9 + 8 - 13)\sqrt{2} \\
        &= 4\sqrt{2}
        \end{aligned}
        $
    \end{itemize}

    \item \begin{itemize}
        \item \textbf{Calcul de $c^2$ :} \\
        $
        \begin{aligned}
        c^2 &= (\sqrt{2} - 3)^2 \\
        &= (\sqrt{2})^2 - 2 \times \sqrt{2} \times 3 + 3^2 \\
        &= 2 - 6\sqrt{2} + 9 \\
        &= 11 - 6\sqrt{2}
        \end{aligned}
        $

        \item \textbf{Déduisons-en que $p$ est un rationnel :} \\
        \textit{(Remarque : dans l'expression de $p$, le terme $3\sqrt{5}$ au dénominateur est une coquille classique d'énoncé pour $11$, afin de faire apparaître $c^2 = 11 - 6\sqrt{2}$. Voici la résolution avec $11 - 6\sqrt{2}$ au dénominateur :)} \\[0.2cm]
        $
        \begin{aligned}
        p &= \dfrac{6 - \sqrt{8}}{11 - 6\sqrt{2}} \\[0.2cm]
        &= \dfrac{6 - \sqrt{2^2 \times 2}}{(\sqrt{2} - 3)^2} \\[0.2cm]
        &= \dfrac{6 - 2\sqrt{2}}{(3 - \sqrt{2})^2} \\[0.2cm]
        &= \dfrac{2(3 - \sqrt{2})}{(3 - \sqrt{2})^2} \\[0.2cm]
        &= \dfrac{2}{3 - \sqrt{2}} \\[0.2cm]
        &= \dfrac{2(3 + \sqrt{2})}{(3 - \sqrt{2})(3 + \sqrt{2})} \\[0.2cm]
        &= \dfrac{6 + 2\sqrt{2}}{3^2 - (\sqrt{2})^2} \\[0.2cm]
        &= \dfrac{6 + 2\sqrt{2}}{9 - 2} \\[0.2cm]
        &= \dfrac{6 + 2\sqrt{2}}{7}
        \end{aligned}
        $ \\[0.2cm]
        \textit{(Note : si l'énoncé original était $p = \dfrac{6 - \sqrt{8}}{\sqrt{11 - 6\sqrt{2}}}$ avec un radical au dénominateur, on obtient alors un rationnel exact :)} \\[0.2cm]
        $
        \begin{aligned}
        p &= \dfrac{6 - \sqrt{8}}{\sqrt{11 - 6\sqrt{2}}} \\[0.2cm]
        &= \dfrac{6 - 2\sqrt{2}}{\sqrt{(\sqrt{2} - 3)^2}} \\[0.2cm]
        &= \dfrac{2(3 - \sqrt{2})}{|\sqrt{2} - 3|} \\[0.2cm]
        &= \dfrac{2(3 - \sqrt{2})}{3 - \sqrt{2}} \quad (\text{car } \sqrt{2} < 3 \text{ donc } \sqrt{2} - 3 < 0) \\[0.2cm]
        &= 2
        \end{aligned}
        $ \\
        Comme $2 \in \mathbb{Q}$, $p$ est bien un nombre rationnel égal à $2$.
    \end{itemize}
\end{enumerate}
\end{comment}

\fbox{\textbf{\exo}}

\begin{enumerate}
    \item \textbf{Simplification de $A$ :} \\
    $
    \begin{aligned}
    A &= \dfrac{\sqrt{3 - \sqrt{3}} \times \sqrt{3 + \sqrt{3}}}{\sqrt{6}} \\[0.2cm]
    &= \dfrac{\sqrt{(3 - \sqrt{3})(3 + \sqrt{3})}}{\sqrt{6}} \\[0.2cm]
    &= \dfrac{\sqrt{3^2 - (\sqrt{3})^2}}{\sqrt{6}} \\[0.2cm]
    &= \dfrac{\sqrt{9 - 3}}{\sqrt{6}} \\[0.2cm]
    &= \dfrac{\sqrt{6}}{\sqrt{6}} \\[0.2cm]
    &= 1
    \end{aligned}
    $

    \item \textbf{Simplification de $B$ :} \\
    $
    \begin{aligned}
    B &= \dfrac{-3(\sqrt{5} - \sqrt{2})}{\sqrt{45} - \sqrt{18}} \\[0.2cm]
    &= \dfrac{-3(\sqrt{5} - \sqrt{2})}{\sqrt{3 \times 3 \times 5} - \sqrt{3 \times 3 \times 2}} \\[0.2cm]
    &= \dfrac{-3(\sqrt{5} - \sqrt{2})}{\sqrt{3^2 \times 5} - \sqrt{3^2 \times 2}} \\[0.2cm]
    &= \dfrac{-3(\sqrt{5} - \sqrt{2})}{3\sqrt{5} - 3\sqrt{2}} \\[0.2cm]
    &= \dfrac{-3(\sqrt{5} - \sqrt{2})}{3(\sqrt{5} - \sqrt{2})} \\[0.2cm]
    &= -1
    \end{aligned}
    $

    \item \textbf{Simplification de $C$ :} \\
    $
    \begin{aligned}
    C &= \dfrac{\sqrt{3} - \sqrt{27} + \sqrt{12}}{\sqrt{54}} \\[0.2cm]
    &= \dfrac{\sqrt{3} - \sqrt{3 \times 3 \times 3} + \sqrt{2 \times 2 \times 3}}{\sqrt{3 \times 3 \times 6}} \\[0.2cm]
    &= \dfrac{\sqrt{3} - \sqrt{3^2 \times 3} + \sqrt{2^2 \times 3}}{\sqrt{3^2 \times 6}} \\[0.2cm]
    &= \dfrac{\sqrt{3} - 3\sqrt{3} + 2\sqrt{3}}{3\sqrt{6}} \\[0.2cm]
    &= \dfrac{(1 - 3 + 2)\sqrt{3}}{3\sqrt{6}} \\[0.2cm]
    &= \dfrac{0\sqrt{3}}{3\sqrt{6}} \\[0.2cm]
    &= 0
    \end{aligned}
    $

    \item \textbf{Simplification de $D$ :} \\
    $
    \begin{aligned}
    D &= \dfrac{2 - \sqrt{5}}{2 + \sqrt{5}} - \dfrac{2 + \sqrt{5}}{2 - \sqrt{5}} \\[0.2cm]
    &= \dfrac{(2 - \sqrt{5})^2}{(2 + \sqrt{5})(2 - \sqrt{5})} - \dfrac{(2 + \sqrt{5})^2}{(2 - \sqrt{5})(2 + \sqrt{5})} \\[0.2cm]
    &= \dfrac{\left[2^2 - 2 \times 2 \times \sqrt{5} + (\sqrt{5})^2\right] - \left[2^2 + 2 \times 2 \times \sqrt{5} + (\sqrt{5})^2\right]}{2^2 - (\sqrt{5})^2} \\[0.2cm]
    &= \dfrac{(4 - 4\sqrt{5} + 5) - (4 + 4\sqrt{5} + 5)}{4 - 5} \\[0.2cm]
    &= \dfrac{(9 - 4\sqrt{5}) - (9 + 4\sqrt{5})}{-1} \\[0.2cm]
    &= \dfrac{9 - 4\sqrt{5} - 9 - 4\sqrt{5}}{-1} \\[0.2cm]
    &= \dfrac{-8\sqrt{5}}{-1} \\[0.2cm]
    &= 8\sqrt{5}
    \end{aligned}
    $

    \item \textbf{Simplification de $E$ :} \\
    \begin{comment}
    Simplifions progressivement en partant du radical le plus intérieur : \\
    $
    \begin{aligned}
    E &= 2 \sqrt{37 - \sqrt{\dfrac{21}{25} + \dfrac{1}{25} \times \sqrt{6 + \sqrt{103 - 2\sqrt{\dfrac{9}{4}}}}}} \\[0.2cm]
    &= 2 \sqrt{37 - \sqrt{\dfrac{21}{25} + \dfrac{1}{25} \times \sqrt{6 + \sqrt{103 - 2 \times \dfrac{3}{2}}}}} \\[0.2cm]
    &= 2 \sqrt{37 - \sqrt{\dfrac{21}{25} + \dfrac{1}{25} \times \sqrt{6 + \sqrt{103 - 3}}}} \\[0.2cm]
    &= 2 \sqrt{37 - \sqrt{\dfrac{21}{25} + \dfrac{1}{25} \times \sqrt{6 + \sqrt{100}}}} \\[0.2cm]
    &= 2 \sqrt{37 - \sqrt{\dfrac{21}{25} + \dfrac{1}{25} \times \sqrt{6 + 10}}} \\[0.2cm]
    &= 2 \sqrt{37 - \sqrt{\dfrac{21}{25} + \dfrac{1}{25} \times \sqrt{16}}} \\[0.2cm]
    &= 2 \sqrt{37 - \sqrt{\dfrac{21}{25} + \dfrac{1}{25} \times 4}} \\[0.2cm]
    &= 2 \sqrt{37 - \sqrt{\dfrac{21}{25} + \dfrac{4}{25}}} \\[0.2cm]
    &= 2 \sqrt{37 - \sqrt{\dfrac{25}{25}}} \\[0.2cm]
    &= 2 \sqrt{37 - \sqrt{1}} \\[0.2cm]
    &= 2 \sqrt{37 - 1} \\[0.2cm]
    &= 2 \sqrt{36} \\[0.2cm]
    &= 2 \times 6 \\[0.2cm]
    &= 12
    \end{aligned}
    $
    \end{comment}
\end{enumerate}

\end{document}